\section{Introduction}
Since the dawn of computing the weakest link in most information systems has been the people creating and using them. We choose easy to remember password that might be prone to dictionary attacks, then we overcomplicate them just to forget them, so to only have to remember a couple passwords we reuse them across sites, also we easily fall for social engineering attacks, like phishing. Sufficiently long, unique, and complex passwords with proper rate limiting remain impractical to brute-force, so compromises overwhelmingly exploit human factors and credential leaks rather than cryptographic weakness. 

To mitigate the weaknesses of a password only environment one-time codes (SMS/TOTP) ere introduced, this is 2 or Multi Factor Authentication, which is meant to reduce the risk of falling for credential stuffing attacks and make phishing more complicated. 2FA doesn't address the core issue of legacy authentication though.

FIDO2/WebAuthn addresses the limitations at last by replacing shared secrets with origin-bound public-key cryptography and local user verification (biometrics or a device PIN). Private keys never leave the authenticator, and the server stores only public keys, eliminating password reuse and greatly reducing the value of stolen databases. Because assertions are tied to the requesting origin, generic phishing and OTP-relay attacks fail by design. In practice, passkeys (discoverable credentials) further streamline the flow across platforms using platform sync or roaming security keys.

This paper adopts that progression from passwords to 2FA to FIDO2 
In the coming section we will walk through 
\begin{itemize}
    \item Mapping common vulnerabilities
    \item Implementing a fully passwordless FIDO2 system
    \item Creating a plan for a STRIDE based analysis
    \item Testing the systems for vulnerabilities
    \item Comparing the systems in terms of performance and security.
\end{itemize}